\newpage
\Article{61}
We extend to the arithmetic of complex numbers what was said in Article~2 of the preceding paper about the notion of a biquadratic residue and nonresidue. Here, as there, we restrict the investigation to moduli that are prime numbers. It will also usually be tacitly understood that the modulus is chosen to be primary among its associates, that is, $\equiv1$ modulo $2+2i$, and that the numbers whose character, as biquadratic residues or nonresidues, is under consideration are not divisible by the modulus.

For a given modulus, therefore, the numbers not divisible by it could be divided into three classes: the first containing biquadratic residues; the second, those biquadratic nonresidues that are quadratic residues; and the third, quadratic nonresidues.

\SourceLeaf{141}%
But here too it is better to establish two classes in place of the third, so that there are four in all.

With any primitive root taken as the base, biquadratic residues will have indices divisible by $4$, or of the form $4n$; those nonresidues that are quadratic residues will have indices of the form $4n+2$; finally, the indices of quadratic nonresidues will be partly of the form $4n+1$, partly of the form $4n+3$. Four classes do arise in this way, but the distinction between the latter two would not be absolute: it would depend on the choice of the primitive root taken as the base. For it is easy to see that one half of the primitive roots gives a given quadratic nonresidue an index of the form $4n+1$, and the other half an index of the form $4n+3$. To remove this ambiguity, we shall always suppose that a primitive root is chosen for which the index $\frac14(p-1)$ belongs to the number $+i$ (compare Article~55, VI). This gives a classification that we can state more concisely, independently of primitive roots, as follows.

\Needspace{10\baselineskip}
Let the \emph{first} class contain those numbers $k$ for which $k^{\frac14(p-1)}\equiv1$; these numbers are biquadratic residues of the modulus.

Let the \emph{second} class contain those for which $k^{\frac14(p-1)}\equiv i$.

The \emph{third}, those for which $k^{\frac14(p-1)}\equiv-1$.

Finally, the \emph{fourth}, those for which $k^{\frac14(p-1)}\equiv-i$.

The third class will contain those biquadratic nonresidues that are quadratic residues; quadratic nonresidues will be distributed between the second and fourth classes.

We shall assign to the numbers in these classes the \emph{biquadratic characters} $0,1,2,3$, respectively. If we define the character $\lambda$ of a number $k$ modulo $m$ as the exponent of that power of $i$ to which $k^{\frac14(p-1)}$ is congruent, then evidently characters congruent modulo $4$ must be regarded as equivalent. For the time being, this notion is restricted to moduli that are prime numbers: in a continuation of these investigations we shall show how it can also be applied to composite moduli.

\Article{62}
To make an extensive examination of examples concerning the characters of numbers easier, we append a compact table here. With its help, the character of any given number with respect to a modulus whose norm does not exceed $157$ is obtained with little work, provided that the following observations are taken into account.

\SourceLeaf{142}%
Since the character of a composite number is equal to, or congruent modulo $4$ to, the sum of the characters of its individual factors, it suffices if we can assign the characters of prime numbers for a given modulus. Further, since the characters of the units $-1,i,-i$ are evidently congruent to the numbers $\frac12(p-1),\frac14(p-1),\frac34(p-1)$ modulo $4$, it will also suffice to give the characters of numbers that are primary among their associates. Finally, since moduli congruent modulo $m$ have the same character, it suffices to enter in the table the characters of numbers contained in a system of absolutely least residues. In addition, reasoning similar to that in Article~58 shows that, if for the modulus $a+bi$ the character of $A+Bi$ is $\lambda$, while for the modulus $a-bi$ $\lambda'$ is the character of $A-Bi$, then we always have $\lambda\equiv-\lambda'\pmod4$, or, equivalently, $\lambda+\lambda'$ is divisible by $4$. It therefore suffices to enter in the table moduli in which $b$ is either $0$ or positive.

\Needspace{16\baselineskip}
Thus, for example, if the character of $11-6i$ with respect to the modulus $-5-6i$ is sought, we substitute the numbers $11+6i,-5+6i$ for them. We then determine (Article~43) the absolutely least residue of $11+6i$ modulo $-5+6i$, which is $-1-4i=-1\times(1+4i)$. Since for the modulus $-5+6i$ the character of $-1$ is $30$, while the character of $1+4i$, from the table, is $2$, we obtain $32$, or $0$, as the character of $11+6i$ for the modulus $-5+6i$; and, by the last observation, this is also the character of $11-6i$ for the modulus $-5-6i$. Likewise, if the character of $-5+6i$ with respect to the modulus $11+6i$ is sought, its absolutely least residue $1-5i$ is factored into $-i\;1+i,\;3-2i$, to which the characters $117,0,1$ correspond. Hence the character sought will be $118$, or $2$; the same character must also be assigned to $-5-6i$ with respect to the modulus $11-6i$.

\[
\begin{array}{r|c|l}
\text{Modulus.}&\text{Character.}&\text{Numbers.}\\\hline
-3&3&1+i\\
+3+2i&3&1+i\\
+1+4i&1&-1+2i\\
&3&1+i\\
-5+2i&0&-1-2i\\
&1&1+i\\
&2&-1+2i\\
-1+6i&0&-3\\
&1&1+i,\ -1+2i
\end{array}
\]

\SourceLeaf{143}%
\[
\begin{array}{r|c|l}
\text{Modulus.}&\text{Character.}&\text{Numbers.}\\\hline
-1+6i&2&-1-2i\\
+5+4i&0&1+i\\
&1&-3\\
&3&-1+2i,\ -1-2i\\
-7&0&-3\\
&1&-1+2i,\ 3-2i\\
&2&1+i\\
&3&-1-2i,\ 3+2i\\
+7+2i&0&1+i,\ 3+2i,\ 3-2i,\ 1-4i\\
&1&-3\\
&2&-1-2i,\ 1+4i\\
&3&-1+2i\\
-5+6i&0&1+i,\ -3,\ 3+2i,\ 3-2i\\
&1&1-4i\\
&2&1+4i\\
&3&-1+2i,\ -1-2i\\
-3+8i&0&-1+2i,\ 3-2i,\ 1-4i\\
&1&1+i,\ 3+2i\\
&2&-3\\
&3&-1-2i,\ 1+4i,\ -5+2i\\
+5+8i&0&-1-2i\\
&1&-5-2i,\ -1+6i\\
&2&-1+2i,\ 3-2i\\
&3&1+i,\ -3,\ 3+2i,\ 1+4i,\ 1-4i\\
+9+4i&0&-1+2i,\ 3+2i\\
&1&1+i,\ -1-2i,\ 3-2i\\
&2&-3,\ 1+4i\\
&3&1-4i,\ -5+2i\\
-1+10i&0&1+i,\ -1+2i,\ -1-2i,\ 3+2i\\
&1&-3\\
&2&3-2i,\ -5+2i,\ 5-4i\\
&3&1+4i,\ 1-4i
\end{array}
\]

\SourceLeaf{144}%
\[
\begin{array}{r|c|l}
\text{Modulus.}&\text{Character.}&\text{Numbers.}\\\hline
+3+10i&1&1+i,\ -1-2i,\ 1-4i\\
&2&-3,\ 3+2i,\ 1+4i,\ -5-2i\\
&3&-1+2i,\ 3-2i\\
-7+8i&0&1+i,\ -7\\
&1&3+2i,\ 3-2i,\ 1-4i,\ -5-2i\\
&2&-1-2i,\ 1+4i,\ -5+2i,\ -1-6i\\
&3&-1+2i,\ -3,\ -1+6i\\
-11&0&-3\\
&1&1+i,\ 3-2i,\ 1+4i,\ -5+2i,\ 5+4i\\
&2&-1+2i,\ -1-2i\\
&3&3+2i,\ 1-4i,\ -5-2i,\ 5-4i\\
-11+4i&0&1+i,\ -1+2i,\ 3+2i,\ 5+4i\\
&1&-1-2i,\ -1+6i\\
&2&-5+2i\\
&3&-3,\ 3-2i,\ 1+4i,\ 1-4i,\ -5-2i\\
+7+10i&0&1+4i,\ 1-4i,\ -1+6i,\ -1-6i\\
&1&-1+2i,\ 3+2i,\ -5+2i\\
&2&1+i,\ 3-2i\\
&3&-1-2i,\ -3,\ -5-2i\\
+11+6i&0&1+i,\ -1+2i,\ -3,\ 1+4i,\ 1-4i,\ -7\\
&1&-1-2i,\ 3+2i,\ 3-2i\\
&2&-5-2i,\ -1+6i,\ 5-4i\\
&3&-5+2i,\ 5+4i,\ 7-2i.
\end{array}
\]

\Article{63}
We shall now seek, by examining examples, common criteria for the moduli for which a given prime number has the same character. We always suppose the moduli to be primary among their associates, that is, numbers $a+bi$ for which either $a\equiv1,b\equiv0$ or $a\equiv3,b\equiv2\pmod4$.

\Needspace{15\baselineskip}
For the number $1+i$, with which we begin, the rule suggested by the examples is more easily grasped if we separate moduli of the first kind, for which $a\equiv1,b\equiv0$, from moduli of the second kind, for which $a\equiv3,b\equiv2$. With the help of the table in the preceding article, we find the following correspondences:
\SourceLeaf{145}%
\[
\begin{array}{c|l}
\text{Character}&\text{Moduli of the first kind.}\\\hline
0&5+4i,\ -7+8i,\ -7-8i,\ -11+4i\\
1&1-4i,\ -3+8i,\ -3-8i,\ 9+4i,\ -11\\
2&5-4i,\ -7,\ -11-4i\\
3&-3,\ 1+4i,\ 5+8i,\ 5-8i,\ 9-4i
\end{array}
\]
If we consider these seventeen examples carefully, in all of them we find the character $\equiv\frac14(a-b-1)\pmod4$.

Likewise, the correspondences are
\[
\begin{array}{c|l}
\text{Character}&\text{Moduli of the second kind.}\\\hline
0&3-2i,\ -1-6i,\ 7+2i,\ -5+6i,\ -1+10i,\ 11+6i\\
1&-5+2i,\ -1+6i,\ 7-2i,\ -1-10i,\ 3+10i\\
2&-1+2i,\ -5-2i,\ 3-10i,\ 7+10i\\
3&-1-2i,\ 3+2i,\ -5-6i,\ 7-10i,\ 11-6i
\end{array}
\]
In all these twenty examples, with a little attention, the character is found to be $\equiv\frac14(a-b-5)\pmod4$.

These two rules can easily be combined into one valid for both kinds of moduli, if we consider that $\frac14bb$ is $\equiv0$ for moduli of the first kind and $\equiv1\pmod4$ for those of the second. Thus the character of $1+i$ with respect to any prime modulus that is primary among its associates is $\equiv\frac14(a-b-1-bb)\pmod4$.

We should note in passing that, since $(b+1)^2$ always has the form $8n+1$, or, equivalently, $\frac14(2b+bb)$ is even, this character is always even or odd according as $\frac14(a+b-1)$ is even or odd. This agrees with the rule for the quadratic character stated in Article~58.

Since $\frac14(a-b-1),\frac14(a-b+3)$ are integers, one even and the other odd, their product will be even, or, equivalently, $\frac18(a-b-1)(a-b+3)\equiv0\pmod4$. Hence the following expression can also be used in place of the one given for the biquadratic character:
\[
\frac14(a-b-1-bb)-\frac18(a-b-1)(a-b+3)=\frac18(-aa+2ab-3bb+1)
\]
This form also has the advantage that it is not restricted to primary moduli, but supposes only that $a$ is odd and $b$ even. For under this supposition either $a+bi$ or $-a-bi$ will evidently be primary among its associates, and the value of that formula is the same for either modulus.
